% Part: computability
% Chapter: computability-theory
% Section: russells-paradox

\documentclass[../../../include/open-logic-section]{subfiles}

\begin{document}

\olfileid{cmp}{thy}{rus}
\olsection{રસેલના વિરોધાભાસ સાથે સરખામણી}

આ વિભાગની દલીલોને રસેલના વિરોધાભાસ સાથે સરખાવવાથી વાત વધુ
સ્પષ્ટ થાય છે:
\begin{enumerate}
\item રસેલનો વિરોધાભાસ: $S = \Setabs{x}{x \notin x}$ રાખો. તો
  $S \in S$ ત્યારે અને માત્ર ત્યારે જ્યારે $S \notin S$,
  જે વિરોધાભાસ છે.
  \sourcecorrection{OLCT-003}{સ્થિર અંગ્રેજી મૂળ અહીં
  $X \notin S$ લખે છે; અગાઉ વ્યાખ્યાયિત કરેલો ગણ $S$ છે અને
  સ્વ-સભ્યતાના આ વિરોધાભાસમાં $S \notin S$ જ જોઈએ.}

  \emph{નિષ્કર્ષ:} આવો ગણ~$S$ નથી. ``બધા ગણોનો ગણ'' અસ્તિત્વમાં
  છે એમ માનવું ગણસિદ્ધાંતના બીજા સ્વયંસિદ્ધાંતો સાથે અસંગત છે.

\item રસેલના વિરોધાભાસનું એક રૂપાંતર: બધા વિધેયોના ગણમાંથી
  $\{ 0, 1 \}$માં જતું ``વિધેય'' $F$ નીચે પ્રમાણે વ્યાખ્યાયિત કરો:
  \[
  F(f) =
  \begin{cases}
    1 & \text{જો $f$ પોતે $f$ના પ્રદેશમાં હોય અને $f(f) = 0$} \\
    0 & \text{અન્યથા}
  \end{cases}
  \]
  સમાન દલીલ બતાવે છે કે $F(F) = 0$ ત્યારે અને માત્ર ત્યારે જ્યારે
  $F(F) = 1$, જે વિરોધાભાસ છે.

  \emph{નિષ્કર્ષ:} $F$ વિધેય નથી. ``બધા વિધેયોનો ગણ'' કોઈ વિધેયનો
  પ્રદેશ બની શકે એટલો નાનો નથી.

\item વિકર્ણીકરણ દલીલ: આંશિક સંગણનીય વિધેયોની પરિગણના
  $f_0$, $f_1$, \dots લો અને $G \colon \Nat \to
  \{ 0, 1 \}$ને નીચે પ્રમાણે વ્યાખ્યાયિત કરો:
  \[
  G(x) =
  \begin{cases}
    1 & \text{જો $f_x(x)\downarrow = 0$} \\
    0 & \text{અન્યથા}
  \end{cases}
  \]
  $G$ સંગણનીય હોય, તો કોઈ $k$ માટે તે વિધેય $f_k$ જ હશે.
  પરંતુ પછી $G(k) = 1$ ત્યારે અને માત્ર ત્યારે જ્યારે $G(k) = 0$,
  જે વિરોધાભાસ છે.

  \emph{નિષ્કર્ષ:} $G$ સંગણનીય નથી. નોંધો કે ગણસિદ્ધાંતના
  સ્વયંસિદ્ધાંતો મુજબ $G$ હજી પણ વિધેય છે; અહીં વિરોધાભાસ
  નથી, માત્ર ભેદ વધુ સ્પષ્ટ થાય છે.
\end{enumerate}

આંશિક વિધેયો, સંગણનીય વિધેયો, આંશિક સંગણનીય વિધેયો વગેરેની
ચર્ચામાં ગૂંચવણ થઈ શકે છે. $\Nat$માંથી $\Nat$માં જતાં બધાં
આંશિક વિધેયોનો ગણ ઘણાં પદાર્થો ધરાવે છે. તેમાંના કેટલાંક સંપૂર્ણ
છે, કેટલાંક સંગણનીય છે, કેટલાંક સંપૂર્ણ અને સંગણનીય બંને છે અને
કેટલાંક બેમાંથી એકેય નથી. યાદ રાખો કે માત્ર ``વિધેય'' કહીએ ત્યારે
સામાન્ય રીતે સંપૂર્ણ વિધેયનો અર્થ લઈએ છીએ. આમ, નીચેના પ્રકારો છે:
\begin{enumerate}
\item સંગણનીય વિધેયો
\item સંપૂર્ણ ન હોય એવા આંશિક સંગણનીય વિધેયો
\item સંગણનીય ન હોય એવા વિધેયો
\item સંપૂર્ણ પણ ન હોય અને સંગણનીય પણ ન હોય એવા આંશિક વિધેયો
\end{enumerate}
આ ભેદ સમજવા માટે $\Nat$માંથી $\Nat$માં જતાં બધા આંશિક વિધેયો
દર્શાવતો મોટો ચોરસ દોરો અને તેમાં અનુક્રમે સંપૂર્ણ વિધેયો તથા
સંગણનીય આંશિક વિધેયો દર્શાવતા બે આચ્છાદિત પ્રદેશો ચિહ્નિત કરો.
આ રેખાંકનમાં મળતા દરેક પ્રદેશમાં આવતો એક પદાર્થ વર્ણવી શકો કે
નહીં તે તપાસવું સારી કસરત છે.

\end{document}
